\documentclass{article}
\title{Early Mobilization After Hip Fracture Surgery in Older Adults}
\date{}
\begin{document}
\maketitle

\begin{abstract}
We tested whether mobilization within 24 hours of hip fracture surgery reduces readmission in older adults. In a randomized trial of 480 patients, early mobilization reduced 30-day readmission from 25 percent to 15 percent. A cohort of 310 patients showed the same direction.
\end{abstract}

\noindent Keywords: hip fracture, early mobilization, readmission, geriatric medicine

\section{Introduction}

Older adults readmitted within 30 days of hip fracture surgery face a higher risk of death (Brown, 2018). Early mobilization may lower the risk of readmission, but only two small trials have tested it (Okoye, 2022).

A randomized trial and a cohort study of routine care address that gap.

\section{Trial Methods}

The trial compared early with usual mobilization in two parallel groups.

A total of 480 patients aged 65 years or older were randomized to early or usual mobilization. Allocation was concealed in sealed envelopes.

We defined early mobilization as standing with a physiotherapist within 24 hours of surgery. Patients then walked twice a day. Usual care began on the second day.

The trial followed the Consolidated Standards of Reporting Trials (CONSORT) guideline.

The authors enrolled the patients on the surgical ward.

Patients were eligible if they could walk unaided before the fracture.

The trial had one limitation: patients and physiotherapists could not be blinded to the allocation.

The primary outcome was readmission within 30 days. Secondary outcomes were length of stay and mortality at 90 days. Outcome assessors did not know the allocation.

Readmission was recorded from hospital records for 30 days after discharge.

\section{Trial Results}

Early mobilization reduced 30-day readmission from 25 percent to 15 percent. The relative risk (RR) was 0.60 (95% CI [0.41, 0.87], p = .004). One readmission was avoided for every ten patients mobilized early.

Each group held 240 patients. All 480 patients completed follow-up and entered the analysis. With no loss to follow-up, attrition cannot account for the difference in readmission.

By 30 days, a heart attack had occurred in two patients, both in the usual care group. The numbers were too small to compare the groups.

Mobilization began a median of 11 hours after surgery in the early group. The median lay inside the 24-hour target.

Length of stay had a median of 7 days in each group. Mortality at 90 days was 6 percent after early mobilization and 8 percent after usual care (RR 0.75, 95% CI [0.39, 1.43]). The trial thus gave no evidence that early mobilization changed either secondary outcome.

\section{Cohort Methods}

The cohort drew on the records of 310 consecutive patients aged 65 years or older. Mobilization time was taken from nursing notes. Mobilization within 24 hours of surgery counted as early and any later mobilization as usual.

Records were incomplete for 18 patients, who were excluded. Their exclusion may have biased the cohort toward patients with complete nursing notes.

The ethics committee approved the use of the nursing notes after a hassle that lasted four months.

The primary outcome matched the trial outcome. Each patient's readmissions within 30 days of discharge came from hospital records.

The analysis adjusted for age and the Charlson comorbidity index, a weighted count of chronic diseases.

\section{Cohort Results}

Readmission occurred in 22 percent of the 140 patients mobilized early and in 31 percent of the 170 patients with usual mobilization. After adjustment, the odds ratio for readmission was 0.68 (95% CI [0.47, 0.98]), against a crude odds ratio of 0.63. Age and comorbidity probably explained little of the association.

Three cases of deep vein thrombosis were recorded in the cohort, a rate of 1 percent. All three occurred in patients with usual mobilization. Usual mobilization may have added to the risk of thrombosis.

\section{Discussion}

The study adds randomized and observational evidence to the two small trials of early mobilization (Okoye, 2022).

The trial and the cohort agree in direction, which suggests that early mobilization reduces readmission. The trial is randomized; the cohort is observational. Agreement between the two designs strengthens the finding.

Both studies come from one hospital, so the effect may differ elsewhere.

Healthier patients, who could be mobilized earlier, may account for part of the cohort association. Adjustment for age and comorbidity changed the odds ratio only from 0.63 to 0.68, which limits that explanation.

The effect shows in readmission, not in length of stay.

Hospitals that wish to test early mobilization could schedule physiotherapy on the day of surgery.

\section{Conclusion}

Early mobilization within 24 hours may reduce readmission after hip fracture surgery in older adults. Future trials should test it in at least three hospitals.


\end{document}
